% !PRINTED 448
\weberpage{448}
%<<CONT
\weq{22}{\begin{aligned}
(k)^p &\equiv (kl) & &(\mathrm{mod}\ \mathfrak P),\\
(kl)^{p'} &\equiv (kll') & &(\mathrm{mod}\ \mathfrak P'),\\
(kll')^{p''} &\equiv (kll'l'') & &(\mathrm{mod}\ \mathfrak P'')
\end{aligned}}
\[
\cdot\ \cdot\ \cdot\ \cdot\ \cdot\ \cdot\ \cdot\ \cdot\ \cdot\ \cdot\ \cdot
\]
Let us now assume that the class equation splits, so that for a
variable $t$
\[
H_m(t) = H_1(t) H_2(t),
\]
then, since $k, k'$ were arbitrary classes, we may choose $(k)$ among
the roots of $H_1(t)$ and $(k')$ among those of $H_2(t)$. In the
chain (22) the initial term $(k)$ therefore belongs to the roots $(k_1)$
of $H_1(t)$, and the final term $(k')$ to the roots $(k_2)$ of $H_2(t)$.
Since the passage from the roots of the one factor to those of
the other must occur at least once in the chain (22),
there is a pair of roots $(k_1), (k_2)$ which, for some
prime $p$ and a prime ideal $\mathfrak P$ dividing it, satisfies the
congruence
\[
(k_1)^p \equiv (k_2) \quad (\mathrm{mod}\ \mathfrak P)
\]
Now since $H_1(k_1) = 0$, it follows by the inference often applied that
\[
H_1(k_1)^p \equiv H_1(k_1^p) \equiv H_1(k_2) \equiv 0 \quad (\mathrm{mod}\ \mathfrak P).
\]
Since now
\[
H_1(k_2) = \prod^{k_1} [(k_2) - (k_1)]
\]
one of the differences $(k_2) - (k_1)$ must be divisible by $\mathfrak P$,
which is not possible, since $p$ does not divide the discriminant
of $H$.

\gesperrt{With this the irreducibility of the class equation is proved; the degree of the field $\mathfrak K$ is established as equal
to twice the class number $h$, and
the group $\mathfrak S$ found in the preceding section
is recognised as the true group of the class equation.}

8. Accordingly we can decompose all primes $p$, \gesperrt{apart
from a finite number of exceptions} (those which divide $2D$
or the discriminant of $H$) into their prime factors
in the field $\mathfrak K$.

If $(D, p) = +1$, then $p$ splits and field $\Omega$ % [sic] German print has "und Körper"; grammatically and mathematically expected "im Körper"
into two conjugate prime ideals of the first degree. If $p$ is of index $\varepsilon$, then
every prime ideal $\mathfrak p$ dividing $p$ is of degree $\varepsilon$. If therefore
\[
p = \mathfrak P_1 \mathfrak P_2 \ldots \mathfrak P_r,
\]
%CONT>>
